\subsection{\texorpdfstring{\(T^{n}\) 的子群与商群}{T 的 n 次幂的子群与商群}}

让我们把 \(\mathbf{T}^n\) 与 \(\mathbf{R}^n/\mathbf{Z}^n\) 等同起来，并设 \(\varphi\) 是 \(\mathbf{R}^n\) 到 \(\mathbf{R}^n/\mathbf{Z}^n\) 上的典范同态。 \(\mathbf{T}^n\) 的每个子群都形如 \(\mathbf{G} = \varphi(\mathbf{H})\) ，其中 \(\mathbf{H}\) 是 \(\mathbf{R}^n\) 的包含 \(\mathbf{Z}^n\) 的子群，且它同构于 \(\mathbf{H}/\mathbf{Z}^n\) （第 III 章 § 2 第 7 小节命题 20）； \(\mathbf{G}\) 在 \(\mathbf{T}^n\) 中为闭的必要充分条件是 \(\mathbf{H}\) 在 \(\mathbf{R}^n\) 中为闭（第 I 章 § 3 第 4 小节）。于是，为了求出 \(\mathbf{T}^n\) 的闭子群，我们必须确定一切闭子群 \(\mathbf{H}\) （ \(\mathbf{R}^n\) 的，包含 \(\mathbf{Z}^n\) ）；为此我们将利用第 3 小节的命题 6，并且先来确定子群 \(\mathbf{H}^*\) （与 \(\mathbf{H}\) 相伴的）。由于 \(\mathbf{Z}^n\) 与其自身相伴，我们有 \(\mathbf{H}^* \subset \mathbf{Z}^n\) ；因而（第 1 小节）存在基 \((\boldsymbol{a}_i)_{1 \leqslant i \leqslant n}\) （ \(\mathbf{R}^n\) 的，它生成 \(\mathbf{Z}^n\) ），以及 \(\mathbf{H}^*\) 的（关于环 \(\mathbf{Z}\) 的）由 \(p\) 个点 \(\boldsymbol{b}_i\) （ \(1 \leqslant i \leqslant p\) ）组成的基，使得 \(\boldsymbol{b}_i = e_i \boldsymbol{a}_i\) （ \(1 \leqslant i \leqslant p\) ），其中诸 \(e_i\) 是满足同余关系 \(e_{i+1} \equiv o (\text{mod } e_i)\) （对 \(1 \leqslant i \leqslant p - 1\) ）的整数。设 \((\boldsymbol{a}_i')\) 是与 \((\boldsymbol{a}_i)\) 对偶的基；则 \(\boldsymbol{u} = \sum_{i=1}^{n} u_i \boldsymbol{a}_i'\) 属于 \((\mathrm{H}^*)^* = \mathrm{H}\) ，当且仅当 \(u_i e_i \in \mathbf{Z}\) （对 \(1 \leqslant i \leqslant p\) ）；换言之， \(\mathbf{H}\) 是由 \(\boldsymbol{a}_{p+1}'\) ， \ldots, \(\boldsymbol{a}_n'\) 所生成的向量子空间 V 与由 \(p\) 个点

\[
e _ {i} ^ {- 1} \boldsymbol {a} _ {i} ^ {\prime} \quad (\mathrm{I} \leqslant i \leqslant p).
\]

另一方面， \(\mathbf{Z}^n\) 是 \(\mathbf{V} \cap \mathbf{Z}^n\) 与 \(\mathbf{K} \cap \mathbf{Z}^n\) 的直和，因为诸 \(a_i'\) （ \(i \leqslant i \leqslant n\) ）生成 \(\mathbf{Z}^n\) 。于是商群 \(\mathrm{H} / \mathbf{Z}^n\) 同构于 \((\mathrm{V} / (\mathrm{V} \cap \mathbf{Z}^n)) \times (\mathrm{K} / (\mathrm{K} \cap \mathbf{Z}^n))\) （第 III 章 § 2 第 9 小节命题 26 的推论）； \(\mathrm{V} / (\mathrm{V} \cap \mathbf{Z}^n)\) 同构于 \(\mathbf{T}^{n-p}\) ，而 \(\mathbf{K} / (\mathbf{K} \cap \mathbf{Z}^n)\) 是一个有限群，是 \(p\) 个循环群的直和，这些循环群的阶分别为 \(e_i\) （ \(i \leqslant i \leqslant p\) ）（参见第 1 小节）。

沿用同样的记号， \(T^{n}\) 的每个豪斯多夫商群都形如 \(\mathbf{T}^{n}/\varphi(\mathbf{H})\) ，且同构于 \(R^{n}/H\) （第 III 章 § 2 第 7 小节命题 22 的推论）；若 W 是由 K 生成的向量子空间，则 \(R^{n}/H\) 同构于 W/K（第 III 章 § 2 第 9 小节命题 26 的推论），即同构于 \(T^{p}\) 。综上所述：

命题 11. \(\mathbf{T}^n\) 的每个闭子群都同构于形如 \(\mathbf{T}^h \times \mathbf{F}\) 的群（ \(0 \leqslant h \leqslant n\) ），其中 \(\mathbf{F}\) 是一个有限交换群，且使得 \(\mathbf{F}\) 是其直和的循环子群的最小个数 \(\leqslant n - h\) 。 \(\mathbf{T}^n\) 的每个豪斯多夫商群都同构于形如 \(\mathbf{T}^h\) 的群（ \(0 \leqslant h \leqslant n\) ）。

特别地 \((n = 1)\) ：

推论. T 的每个闭子群，除 T 自身外，都是有限循环群。T 的每个豪斯多夫商群，除 \(\{o\}\) 外，都同构于 T。

